\chapter{Introduction and Motivation}
% ===================================================================

Physics is, at its core, a theory of cause and effect: given the present state
of a system and the laws governing its evolution, what future states are
possible, and with what probabilities?  This paper argues that
\emph{graph-structured lambda theories} (GSLTs) --- the natural categorical
setting for process calculi and rewriting systems --- provide an unusually
transparent foundation for this question.

The central thesis is that the structures of classical and quantum mechanics
(Lagrangian, Hamiltonian, path integral, conservation laws, time symmetry) are
not additional assumptions layered on top of a dynamical theory; rather, they
\emph{emerge from} the combinatorial structure of a GSLT once that structure is
equipped with modest additional data.  The required data amounts to:

\begin{itemize}[itemsep=2pt]
  \item a \emph{weight map} assigning complex amplitudes to rewrite steps,
        indexed by a context-decorated Hennessy--Milner logic;
  \item a \emph{cost map} assigning resource consumption to rewrite steps,
        indexed by the same logic;
  \item a \emph{vectorial account} tracking available resources across
        parallel components of a term.
\end{itemize}

Everything else --- the path integral, interference, conservation of energy,
time-reversal symmetry --- follows by construction.

\section{A First Glimpse: The Rosetta Stone}

The following table is a preview of the correspondence developed in this paper.
Each row is a theorem or construction, not a definition by analogy.
The reader may find it useful as a map to consult while reading;
a fully annotated version appears in Section~\ref{sec:synthesis}.

\begin{center}
\renewcommand{\arraystretch}{1.4}
\begin{tabular}{>{\itshape}p{4.8cm} p{6.8cm}}
\toprule
\normalfont\textbf{GSLT} & \textbf{Physics} \\
\midrule
Term                          & State / initial condition \\
Rewrite rule                  & Law of motion \\
Rewrite path                  & Trajectory \\
Synchronization tree          & Causal graph \\
Minimum path length           & Proper time \\
Equations $\eqs$              & Gauge equivalence \\
Reversible envelope $\GSLT^\dagger$ & CPT-symmetric theory \\
Weight map (complex-valued)   & Lagrangian / path amplitude $e^{iS/\hbar}$ \\
Cost map + vectorial account  & Hamiltonian + conserved charges \\
Conservation (from reversibility) & Conservation of energy \\
Symmetry of cost map          & Noether current \\
Complex weight map            & Quantum probability amplitude \\
Sum over paths                & Feynman path integral \\
\bottomrule
\end{tabular}
\end{center}

\section{Running Example: The Rho Calculus}

Throughout this paper we use the Rho calculus~\cite{meredith2005rho} as a
running example because it exhibits, in compact form, many of the features of
interest: reflection (processes and names are mutually definable), parallel
composition, and a synchronization mechanism that serves as the prototype for
the causal graphs we study.

The grammar of the Rho calculus is:
\begin{align}
  P, Q &\;::=\; 0 \;\mid\; \mathtt{for}(y \leftarrow x)\,P \;\mid\;
               x!(Q) \;\mid\; P \mid Q \;\mid\; {*}x \\
  x, y &\;::=\; @P
\end{align}
where $y$ in $\mathtt{for}(y \leftarrow x)\,P$ is a bound variable.
Structural equivalence $\equiv$ is the smallest equivalence relation containing
$\alpha$-equivalence making $(P,\,|\,,0)$ a commutative monoid.
The single reduction rule is:
\[
  \mathtt{for}(y \leftarrow x)\,P \;\mid\; x!(Q)
  \;\rewrite\;
  P\{@Q / y\}
\]

\section{Organization}

Section~\ref{sec:gslt} defines GSLTs and their category.
Section~\ref{sec:reversibility} constructs the time-symmetric envelope $\GSLT^\dagger$.
Section~\ref{sec:causal} identifies the synchronization tree with a causal graph.
Section~\ref{sec:lts-hml} reviews the Milner--Sewell--Leifer construction and
the resulting context-decorated HML.
Section~\ref{sec:weighted} introduces the weight map and the Lagrangian structure.
Section~\ref{sec:cost} introduces the cost map and the Hamiltonian structure.
Section~\ref{sec:extended-modal} extends the modal operators with full dynamical state.
Section~\ref{sec:quantum} lifts weights to complex amplitudes and connects to the
quantum Gillespie algorithm.
Section~\ref{sec:synthesis} collects the structures into a unified picture and
states the main conservation theorem.

% ===================================================================
